\documentclass[10pt,a4paper]{amsart}
\usepackage[margin=19mm]{geometry}
\usepackage{amsmath,amssymb,mathtools}
\usepackage[scr=rsfs,cal=euler]{mathalfa}
\usepackage{xfrac}
\usepackage[unicode,hidelinks,bookmarks=false]{hyperref}
\newcommand{\ie}{i.e.\ }
\newcommand{\cf}{cf.\ }
\newcommand{\resp}{respectively\ }
\newcommand{\Subsection}{\subsection}
\providecommand{\nobreakdash}{\nobreak\hbox{-}}
\setlength{\emergencystretch}{2em}
\multlinegap0pt
\usepackage{amssymb}
\usepackage{mathtools}
\usepackage[shortlabels]{enumitem}
\newtheorem{lemma}{Lemma}
\newtheorem{proposition}{Proposition}
\newcommand{\ip}[2]{\left\langle #1,#2\right\rangle}
\title{Dynamic Averaging on Regular Graphs}
\author{Dean Kraizberg}
\date{Source: arXiv:2607.00966v3 (CC BY 4.0)}
\begin{document}
\maketitle

\noindent\emph{Source and licence.} Dynamic Averaging on Regular Graphs. Version arXiv:2607.00966v3 (CC BY 4.0), distributed by arXiv under the Creative Commons Attribution 4.0 International licence, http://creativecommons.org/licenses/by/4.0/, as recorded in the arXiv licence metadata for this exact version and shown on the abstract page https://arxiv.org/abs/2607.00966v3. Full licence notice: \texttt{licenses/arxiv-2607.00966-CC-BY-4.0.txt}.\par\smallskip

\noindent\textit{Editorial scope.} Counted unit: the article's proposition prop:one-step (source0.tex lines 853--867). The article's complete original proof of it is delivered with it (source0.tex lines 869--1099, one \texttt{proof} environment of 231 lines). No mathematics is written, repaired or re-derived: the statement and the proof are the article's own lines, in the article's own notation, and no retained line is substituted or re-flowed.  Retained uncounted as the source-local context the proof needs: lines 789--798.  Nothing was omitted from any retained span, so the package carries no omission record.  No retained line contains a citation, so the wrapper carries no bibliography and no bibliography entry.  No retained line contains a figure environment, an image-inclusion command or drawing code, so no image is delivered and the package declares no asset dependency.  Editorial formatting only, in addition: an AMS \texttt{amsart} preamble that restates the article's own declarations and macros used by the retained lines, namely the environments \texttt{lemma}, \texttt{proposition} on separate counters, together with the standard packages those lines need.  The retained environments print in this document as Lemma 1, Proposition 1 in order of appearance, whereas the article prints the same items with the numbering of its own full document.  The complete native source of the article as distributed in the arXiv e-print for this exact version is bundled unchanged as \texttt{sources/204200/source0.tex}.
\begin{lemma}
\label{lem:voltage-bound}
Let \(\theta\in \mathbf 1^{\perp}\), and set \(\psi=\mathcal L^+\theta.\)
Then, for every edge \(\{u,v\}\in E\),
\[
|\psi_u-\psi_v|
\le
\frac{\|\theta\|_1}{2}R_{\mathrm{eff}}(u,v)
\]
\end{lemma}

\begin{proposition}
\label{prop:one-step}
Let \(G\) be a connected \(d\)-regular graph. For every zero-sum vector \(\theta\in\ \mathbf 1 ^\perp \) satisfying \(\|\theta\|_1 \le \frac{1}{3(\Re_G+1)}\)
and every \(w\in[0,1]\),
\[
\frac1{|E|}
\sum_{e\in E}
\exp\left(
 w\ip{\theta}{c_e}
+
4\bigl(R(A_e\theta)-R(\theta)\bigr)
\right)
\le1.
\]
\end{proposition}

\begin{proof}
We again denote \(
\psi=\mathcal L^+\theta.\)
Fix an orientation of every edge \(e=\{u,v\}\) and define
\[
\delta_e=\theta_u-\theta_v,
\qquad
z_e=\ip{\theta}{c_e}, \qquad q_e
=
R(A_e\theta)-R(\theta).
\]
Since \(\theta\) has zero sum,
\[
z_e
=
\frac{\theta_u+\theta_v}{2}.
\]
We compute now \(q_e\). Since
\(
A_e\theta
=
\theta-\frac{\delta_e}{2}(e_u-e_v),
\)
the effective resistance part of \(R\) changes by
\[
\begin{aligned}
&\frac d2
\left(
\ip{A_e\theta}{\mathcal L^+A_e\theta}
-
\ip{\theta}{\mathcal L^+\theta}
\right)\\
&\qquad=
-\frac d2\delta_e
\ip{e_u-e_v}{\mathcal L^+\theta}
+
\frac d8\delta_e^2
\ip{e_u-e_v}{\mathcal L^+(e_u-e_v)}\\
&\qquad=
-\frac d2\delta_e(\psi_u-\psi_v)
+
\frac d8\delta_e^2R_{\mathrm{eff}}(u,v).
\end{aligned}
\]
On the other hand,
\[
\begin{aligned}
\|A_e\theta\|_2^2-\|\theta\|_2^2
&=
\left\|
\theta-\frac{\delta_e}{2}(e_u-e_v)
\right\|_2^2
-
\|\theta\|_2^2\\
&=
-\delta_e\ip{\theta}{e_u-e_v}
+
\frac{\delta_e^2}{4}\|e_u-e_v\|_2^2\\
&=
-\delta_e^2+
\frac12\delta_e^2\\
&=
-\frac12\delta_e^2.
\end{aligned}
\]
Multiplying by the coefficient \((\Re_G+2)/4\) in the definition of \(R\), we obtain
\[
q_e
=
-\frac d2\delta_e(\psi_u-\psi_v)
+
\frac{dR_{\mathrm{eff}}(u,v)-(\Re_G+2)}{8}\delta_e^2 
\le
-\frac d2\delta_e(\psi_u-\psi_v)
-
\frac{1}{4}\delta_e^2.
\]
We next sum the first-order terms over all edges. We observe that
\[
\sum_{e\in E}z_e
=
\frac12
\sum_{\{u,v\}\in E}
(\theta_u+\theta_v)
=
\frac d2\sum_{i\in V}\theta_i
=
0
\]
and also,
\[
\sum_{e\in E}z_e^2
=
\frac14
\sum_{\{u,v\}\in E}
(\theta_u+\theta_v)^2
=
\frac d2\|\theta\|_2^2
-
\frac14\ip{\theta}{\mathcal L\theta}.
\]
To calculate the mixed term in \(q_e\), note that
\(
\sum_{e=\{u,v\}\in E} \delta_{e}(\psi_u-\psi_v)=\ip{\theta}{\mathcal L\psi}=\ip{\theta}{\mathcal L \mathcal L^+ \theta}=\|\theta\|_2^2.\) From the previous bound on $q_e$ we obtain
\[
\sum_{e\in E}q_e
\le
\sum_{e=\{u,v\}\in E} -\frac d2  \delta_{e}(\psi_u-\psi_v)
-
\sum_{e\in E} \frac 1 4\delta_e^2
=
-\frac d2\|\theta\|_2^2
-
\frac14\ip{\theta}{\mathcal L\theta}
=
-
\sum_{e\in E}z_e^2
-
\frac12\ip{\theta}{\mathcal L\theta}.
\]
Since $\theta \in \mathbf 1^\perp$, the negative and positive entries of $\theta$ have equal sum, and in particular $\|\theta^-\|_1=\|\theta^+\|_1 = \frac{\|\theta\|_1}{2}$. Thus \(|\delta_e| \le \|\theta\|_1\) for every edge \(e\), and by Lemma~\ref{lem:voltage-bound},
\(
|\psi_u-\psi_v|\le
\frac{\|\theta\|_1}{2}R_{\mathrm{eff}}(u,v)
\le
\frac{\Re_G\|\theta\|_1}{2d}.
\)
Using again the previous calculation of $q_e$ we see
\[
\begin{aligned}
|q_e|
&\le
\frac d2|\delta_e|\,|\psi_u-\psi_v|
+
\frac{(\Re_G+2)-dR_{\mathrm{eff}}(u,v)}{8}\delta_e^2\\
&\le
\frac{\Re_G\|\theta\|_1}{4}|\delta_e|
+
\frac{\Re_G+2}{8}\delta_e^2\\
&\le
\frac{\Re_G\|\theta\|_1}{4}|\delta_e|
+
\frac{\Re_G+2}{8}\|\theta\|_1|\delta_e|\\
&=
\frac{3\Re_G+2}{8}
\|\theta\|_1|\delta_e|.
\end{aligned}
\]
Then
\[
\sum_{e\in E}q_e^2
\le
\left(\frac{3\Re_G+2}{8}\right)^2
\|\theta\|^2_1
\sum_{e\in E}\delta_e^2
=
\left(\frac{3\Re_G+2}{8}\right)^2
\|\theta\|^2_1
\ip{\theta}{\mathcal L\theta}.
\]

We are now ready for the exponential estimate. Denote \(y_e = wz_e+4q_e.\) Notice
\[
\sum_{e\in E}y_e
=
 w\sum_{e\in E}z_e
+
4\sum_{e\in E}q_e
\le
-4\sum_{e\in E}z_e^2
-
2\ip{\theta}{\mathcal L\theta}.
\]
Since \((x+y)^2\le2x^2+2y^2\) we also get,
\[
\sum_{e\in E}y_e^2
\le
2w^2\sum_{e\in E}z_e^2
+
32\sum_{e\in E}q_e^2
\le
2\sum_{e\in E}z_e^2
+32
\left(\frac{3\Re_G+2}{8}\right)^2
\|\theta\|^2_1
\ip{\theta}{\mathcal L\theta}.
\]
Since \(\|\theta\|_1 \le \frac 1{3(\Re_G+1)}\) we have 
\(
32 \left(\frac{3\Re_G+2}{8}\right)^2
\|\theta\|^2_1\le \frac{ (\Re_G+\frac23)^2}{2(\Re_G+1)^2}\le \frac12
\)
and consequently,
\[
\sum_{e\in E}(y_e+y_e^2)
\le
(-4+2)\sum_{e\in E}z_e^2
+
\left(
-2+\frac{1}{2}
\right)
\ip{\theta}{\mathcal L\theta}
\le
0.
\]
Note that
\[
|y_e|
\le
|z_e|+4|q_e|
\le
\|\theta\|_1+\frac{3\Re_G+2}{2} \|\theta\|_1^2
\le
\frac32 \|\theta\|_1\le
\frac12.
\]
For \(|y|\le \frac12\), the inequality \( e^y\le1+y+y^2 \) holds, and by applying it to every \(y_e\) and summing over the edges gives
\[
\begin{aligned}
\sum_{e\in E}e^{y_e}
&\le
\sum_{e\in E}(1+y_e+y_e^2)\\
&=
|E|+
\sum_{e\in E}(y_e+y_e^2)\\
&\le
|E|.
\end{aligned}
\]
Dividing by \(|E|\) proves the proposition.
\end{proof}

\end{document}
